\renewcommand\L{\mathcal{L}}
\newcommand\G{\tilde{\Gamma}}
\begin{problem}[1]
  A three-dimensional spacetime has the line element
  \[%
    \ds^2 = -\left(1 - \frac{2M}{r}\right) \ds t^2 + \left(1 - \frac{2M}{r}\right)^{-1} \dr^2 + r^2 \dd{\phi}^2
  .\]%
  \begin{enumerate}
    \item (10 points) Find the explicit Lagrangian for the variational principle for geodesics in this spacetime in these coordinates.
    \item (10 points) Using the results of (i) write out the components of the geodesic equation by computing them from the Lagrangian.
    \item (5 points) Read off the nonzero Christoffel symbols for this metric from your results in (ii).
  \end{enumerate}
\end{problem}

\begin{solution}[(i)]
  The metric, $g_{\mu\nu}$, is given by:
  \[%
    g_{\mu\nu} = \begin{pmatrix}
      -\left(1 - \frac{2M}{r}\right) & 0 & 0 \\
      0 & \left(1 - \frac{2M}{r}\right)^{-1} & 0 \\
      0 & 0 & r^2
    \end{pmatrix}
  .\]%
  The Lagrangian, $\L$, for the variational principle for geodesics is given by:
  \[%
    \L = \frac{1}{2} g_{\mu\nu} \dot{x}^\mu \dot{x}^\nu
  ,\]%
  where $\dot{x}^\mu = \odv{x^\mu}/{\lambda}$ is a potential term and $\lambda$ is an affine parameter along the geodesic. Thus, we have:
  \[%
    \L = \frac{1}{2} \left[-\left(1 - \frac{2M}{r}\right) \dot{t}^2 + \left(1 - \frac{2M}{r}\right)^{-1} \dot{r}^2 + r^2 \dot{\phi}^2 \right]
  .\qedhere\]%
\end{solution}

\begin{solution}[(ii)]
  The Euler-Lagrange equations for the geodesic motion are given by:
  \[%
    \odv{}{\lambda} \left(\pd{\dot{x}^\mu}\L\right) - \pd{\mu}\L = 0
  .\]%
  We compute these for each coordinate. For $\mu = t$, we have:
  \begin{align*}
    &\quad\odv{}{\lambda} \left(-\left(1 - \frac{2M}{r}\right) \dot{t}\right) = 0 \\
    \implies&\quad \left(1 - \frac{2M}{r}\right) \dot{t} = E
  ,\end{align*}
  where $E$ is a constant of motion. For $\mu = r$, we have:
  \begin{align*}
    &\quad\odv{}{\lambda} \left(\left(1 - \frac{2M}{r}\right)^{-1} \dot{r}\right) - \pd{r}\L = 0 \\
    \implies&\quad \odv{}{\lambda} \left(\left(1 - \frac{2M}{r}\right)^{-1} \dot{r}\right) + \frac{M}{r^2} \dot{t}^2 - \frac{M}{r^2} \left(1 - \frac{2M}{r}\right)^{-2} \dot{r}^2 - r \dot{\phi}^2 = 0
  .\end{align*}
  For $\mu = \phi$, we have:
  \begin{align*}
    &\quad\odv{}{\lambda} \left(r^2 \dot{\phi}\right) = 0 \\
    \implies&\quad r^2 \dot{\phi} = L
  ,\end{align*}
  where $L$ is another constant of motion. Thus, the geodesic equations are:
  \begin{gather*}
    \left(1 - \frac{2M}{r}\right) \dot{t} = E, \quad r^2 \dot{\phi} = L \\
    \odv{}{\lambda} \left(\left(1 - \frac{2M}{r}\right)^{-1} \dot{r}\right) + \frac{M}{r^2} \dot{t}^2 - \frac{M}{r^2} \left(1 - \frac{2M}{r}\right)^{-2} \dot{r}^2 - r \dot{\phi}^2 = 0
  .\qedhere\end{gather*}
\end{solution}

\begin{solution}[(iii)]
  The nonzero Christoffel symbols can be read off from the geodesic equations. We have:
  \begin{gather*}
    \Gamma^t_{tr} = \frac{M}{r^2} \left(1 - \frac{2M}{r}\right)^{-1}, \quad \Gamma_{tt}^r = \frac{M}{r^2} \left(1 - \frac{2M}{r}\right) \\
    \Gamma^r_{rr} = -\frac{M}{r^2} \left(1 - \frac{2M}{r}\right)^{-1}, \quad \Gamma^r_{\phi\phi} = -r \left(1 - \frac{2M}{r}\right) \\
    \Gamma^\phi_{r\phi} = \frac{1}{r}
  .\qedhere\end{gather*}
\end{solution}

\begin{problem}[2 (Christoffel Symbols, aka connection coefficients)]
  Consider a change of basis $\tilde{e}_\mu = A_\mu^\nu e_\nu$.
  \begin{enumerate}
    \item (10 points) Show that the components of the Christoffel symbol in the new basis are related to the components in the old basis by
      \[%
        \G_{\rho\nu}^\mu = A_\rho^\sigma \left(A^{-1}\right)_\tau^\mu \left[A_\tau^\lambda \Gamma_{\sigma\lambda}^\tau + \pd{\sigma}A_\nu^\tau\right]
      .\]%
      Given the above equation, are the Christoffel symbols tensors? Why or why not?

    \item (10 points) Does the difference of two connections, $S_{\rho\nu}^\mu = (\Gamma_1)_{\rho\nu}^\mu - (\Gamma_2)_{\rho\nu}^\mu$ obey the tensor transformation laws?
  \end{enumerate}
\end{problem}

\begin{solution}[(i)]
  The Christoffel symbols are defined by the relation:
  \[%
    \nabla_\mu e_\nu = \Gamma_{\mu\nu}^\rho e_\rho
  .\]%
  Under a change of basis, we have:
  \begin{align}
    \nabla_\mu \tilde{e}_\nu &= \nabla_\mu (A_\nu^\sigma e_\sigma) \tag*{} \\
                             &= \pd{\mu} A_\nu^\sigma e_\sigma + A_\nu^\sigma \nabla_\mu e_\sigma \tag*{} \\
                             &= \pd{\mu} A_\nu^\sigma e_\sigma + A_\nu^\sigma \Gamma_{\mu\sigma}^\tau e_\tau \tag*{} \\
                             &= \left(\pd{\mu} A_\nu^\tau + A_\nu^\sigma \Gamma_{\mu\sigma}^\tau\right) e_\tau \label{eq:christoffel_transform_1}
  .\end{align}
  On the other hand, we also have:
  \begin{equation}\label{eq:christoffel_transform_2}
    \nabla_\mu \tilde{e}_\nu = \G_{\mu\nu}^\rho \tilde{e}_\rho = \G_{\mu\nu}^\rho A_\rho^\tau e_\tau
  \end{equation}
  Equating the right-hand sides of equations~\eqref{eq:christoffel_transform_1} and~\eqref{eq:christoffel_transform_2}, we get:
  \[%
    \G_{\mu\nu}^\rho A_\rho^\tau = \pd{\mu} A_\nu^\tau + A_\nu^\sigma \Gamma_{\mu\sigma}^\tau
  .\]%
  Multiplying both sides by $\left(A^{-1}\right)_\tau^\mu$, we obtain:
  \[%
    \G_{\mu\nu}^\rho = A_\mu^\sigma \left(A^{-1}\right)_\tau^\rho \left[A_\nu^\lambda \Gamma_{\sigma\lambda}^\tau + \pd{\sigma}A_\nu^\tau\right]
  .\]%
  The Christoffel symbols do not transform as tensors because of the presence of the inhomogeneous term $\pd{\sigma}A_\nu^\tau$ in the transformation law. This term prevents the Christoffel symbols from obeying the tensor transformation laws, which require that all terms in the transformation be homogeneous (i.e., involve only the original components and the transformation matrices without derivatives).
\end{solution}

\begin{solution}[(ii)]
  Let $S_{\rho\nu}^\mu = (\Gamma_1)_{\rho\nu}^\mu - (\Gamma_2)_{\rho\nu}^\mu$. Under a change of basis, we have:
  \begin{align*}
    S_{\rho\nu}^\mu \to \G_{\rho\nu}^\mu &= A_\rho^\sigma \left(A^{-1}\right)_\tau^\mu \left[A_\nu^\lambda (\Gamma_1)_{\sigma\lambda}^\tau + \pd{\sigma}A_\nu^\tau\right] - A_\rho^\sigma \left(A^{-1}\right)_\tau^\mu \left[A_\nu^\lambda (\Gamma_2)_{\sigma\lambda}^\tau + \pd{\sigma}A_\nu^\tau\right] \\
                                         &= A_\rho^\sigma \left(A^{-1}\right)_\tau^\mu A_\nu^\lambda \left[(\Gamma_1)_{\sigma\lambda}^\tau - (\Gamma_2)_{\sigma\lambda}^\tau\right] \\
                                         &= A_\rho^\sigma \left(A^{-1}\right)_\tau^\mu A_\nu^\lambda S_{\sigma\lambda}^\tau
  .\end{align*}
  Thus, the difference of two connections, $S_{\rho\nu}^\mu$, transforms as a tensor under a change of basis. This is because the inhomogeneous terms cancel out when taking the difference, leaving only the homogeneous transformation terms that satisfy the tensor transformation laws.
\end{solution}

\begin{problem}[3]
  (15 points) ($\star$) An observer in a circular orbit of circumference $2\pi R$ around a charged, spherical, black hole (Reissner-Nordstr\"om black hole) of mass $M$ and charge $q$ measures local electric and magnetic fields. What are their strengths and orientations? The Reissner-Nordstr\"om metric is given by
  \[%
    \ds^2 = -f(r)\dt^2 + f(r)^{-1}\dr^2 + r^2(\dd{\theta}^2 + \sin^2\theta \dd{\phi}^2)
  ,\]%
  where $f(r) \equiv 1 - 2GM/r + q^2/r$.
\end{problem}

\begin{solution}
  The strengths and orientations of the electric and magnetic fields measured by the observer can be determined using the electromagnetic field tensor, $F_{\mu\nu}$, which is related to the four-potential, $A_\mu$, by:
  \[%
    F_{\mu\nu} = \pd{\mu} A_\nu - \pd{\nu} A_\mu
  .\]%
  For a Reissner-Nordstr\"om black hole, the four-potential is given by:
  \[%
    A_\mu = \left(-\frac{q}{r}, 0, 0, 0\right)
  .\]%
  The nonzero components of the electromagnetic field tensor are:
  \[%
    F_{tr} = -F_{rt} = \pd{t} A_r - \pd{r} A_t = \frac{q}{r^2}
  .\]%
  The electric field, $\vec{E}$, measured by the observer is given by:
  \[%
    E^i = F^{i0} = g^{ij} F_{j0}
  .\]%
  In the observer's local frame, the electric field is:
  \[%
    E^r = g^{rr} F_{rt} = f(r) \frac{q}{r^2}, \quad E^\theta = 0, \quad E^\phi = 0
  .\]%
  The magnetic field, $\vec{B}$, measured by the observer is given by:
  \[%
    B^i = \frac{1}{2} \epsilon^{ijk} F_{jk}
  .\]%
  Since the only nonzero component of $F_{\mu\nu}$ is $F_{tr}$, the magnetic field is zero:
  \[%
    B^r = 0, \quad B^\theta = 0, \quad B^\phi = 0
  .\]%
  Thus, the observer measures an electric field of strength $E^r = f(R) \frac{q}{R^2}$ directed radially outward from the black hole, and no magnetic field.
\end{solution}
